\documentclass[11pt]{letter}
\usepackage[utf8]{inputenc}
\usepackage[T1]{fontenc}
\usepackage[margin=1in]{geometry}
\usepackage{siunitx}
\usepackage{hyperref}

\signature{Jonathan Mace\\for the authors}
\address{Jonathan Mace\\Microsoft Research\\Redmond, WA, USA}
\date{\today}

\begin{document}

\begin{letter}{%
Editor-in-Chief\\
The Journal of the Acoustical Society of America\\
Acoustical Society of America}

\opening{Dear Editor-in-Chief,}

Please consider our manuscript, \emph{``A Constrained-Bubble Model for Gas-Pocket Mechanotransduction in the Human Abdomen: Theoretical Framework and Sensitivity Analysis,''} for publication in \emph{The Journal of the Acoustical Society of America}.

The paper develops a constrained-bubble model for intestinal gas pockets as compliant inclusions embedded in soft tissue. Rather than revisiting whole-cavity resonance, we analyse a distinct local transduction pathway: sub-resonant oscillation of intraluminal gas pockets under airborne infrasound. The model predicts an upper-bound susceptibility range of \SIrange{111}{120}{\decibel} for an illustrative mechanotransduction threshold, giving a \SI{9}{\decibel} spread across physiologically plausible pocket sizes. That range offers a concrete physical explanation for differential vulnerability, including greater susceptibility in individuals with elevated bowel gas such as post-prandial subjects or some patients with irritable bowel syndrome.

We believe the manuscript fits JASA well because it sits at the intersection of biomedical acoustics, bubble dynamics, and wave propagation in soft biological media. The mathematical core is an acoustics problem: pressure-driven oscillation of a compressible inclusion, modified by elastic confinement, radiation loading, and low-frequency pressure equalisation through the gastrointestinal air column. The paper should therefore interest readers working in biomedical ultrasound, bubble acoustics, acoustofluidics, and bioacoustic transduction.

This submission is a companion to our Paper 1 study now under consideration at \emph{Journal of Sound and Vibration}, which showed that whole-abdomen airborne coupling is negligible. The present manuscript is fully self-contained and does not depend on that paper for its theory, data, or conclusions; it addresses a separate mechanism and stands on its own as a biomedical-acoustics contribution.

We also emphasise that the paper treats the mechanism cautiously. The analysis identifies acoustic short-circuiting through the GI air column as the critical uncertainty and presents the sealed-pocket result as an upper bound rather than an established physiological effect. In that respect, the manuscript contributes both a new model and a clear experimental agenda for testing it.

The manuscript is original, not under consideration elsewhere, and approved by both authors. Brian R. Mace's contact email for submission metadata is \texttt{b.mace@auckland.ac.nz}.

\closing{Yours sincerely,}

\end{letter}
\end{document}


